\section{Formative Study}
\label{sec:5}

The formative stage will establish whether the target event occurs, can be labeled with reasonable agreement, and has an acceptable intervention option. Context-triggered studies of proactive interaction provide a precedent for investigating assistance in the situation where it would occur ~\cite{expect,speaker}.

\subsection{Participants and procedure}

\missing{[TO COMPLETE: recruitment, number and size of teams, expertise, inclusion criteria, consent and ethics approval, task protocol, interview prompts, and session duration.]}

\subsection{Annotation and intervention acceptability}

\missing{[TO COMPLETE: event codebook, coder training, inter-rater agreement, disagreements, participant recall procedure, and evidence about acceptable public versus private prompts.]}

\subsection{Technical feasibility}

\missing{[TO COMPLETE: sensor comfort in XR, usable-data coverage, synchronization error, motion artifacts, processing delay, and failed pilot configurations.]}

\subsection{Findings that justify the design}

\missing{[TO COMPLETE with actual evidence only: observed problem frequency, illustrative anonymized episodes, design changes, and implications for D1 to D5. Do not insert hypothetical quotations or pilot statistics.]}
